\paragraph{P-F57, far-increment estimate.}
Use the exact increment $J_n(a)$ of (P-F56ab)--(P-F56ac),
with $n\ge5$, $l=n-1$, $M=2^l$, and $b_n=3^{-1}\pmod M$.
For an odd residue $a\pmod M$, let
\[
 J=\min\{j,M-j\},\qquad j\equiv3a\pmod M,
                 \quad 1\le j<M.
\]
Thus $J$ is an odd integer with $1\le J<M/2$.
To distinguish the integer $J$ from the increment, write
$I_n(a)=J_n(a)$ in this paragraph. Then
\[
 \log_2\frac{I_n(a)}{I_n(b_n)}
       \le 7n^2-\frac{(\log_2J-8)_+^3}{24}.
 \tag{P-F57f1}
\]
In particular, for each fixed $\delta>0$, uniformly for
$J\ge2^{\delta n}$ and all sufficiently large $n$,
\[
       \frac{I_n(a)}{I_n(b_n)}
                    \le2^{-\delta^3n^3/384}.
 \tag{P-F57f2}
\]
For example it suffices here that
$n\ge\max\{5,16/\delta,2688/\delta^3\}$.

\emph{Proof.}
Put $w_i=(-2)^i$ for $0\le i<l$ and $w_l=0$.
We first bound all the active pair factors
\[
       1-\exp\bigl(2\pi i a(1+w_j+w_k)/M\bigr),
                 \qquad 2\le j<k<l.
 \tag{P-F57f3}
\]
The increment includes all of these except $(j,k)=(2,3)$.
Every phase is nonzero modulo one, because its numerator
is odd. Write $\theta=a/M$ and
$y_{jk}=\theta(1+w_j+w_k)$ modulo one.
For an integer $u\in\{0,\ldots,l-1\}$ put
$r=l-u$ and let $N_u$ be the number of pairs in (P-F57f3)
with $\|y_{jk}\|<2^{-u}$.

The map from subsets of $\{0,\ldots,l-1\}$ to residues
modulo $2^l$ that sends a subset to the sum of its $w_i$
is injective. Indeed the difference of two distinct subset
sums has a unique least index $i$ with nonzero coefficient
$1$ or $-1$. Its two-adic valuation is exactly $i<l$.
This also proves injectivity for the pair sums used below;
the zero-weight variable $w_l$ is not included here.

Call an active index early if $j<r$, and late otherwise.
The number of early indices is $e=(r-2)_+$.
There are at most $\binom e2$ early--early pairs.
For each fixed early index $j$, all mixed pair numerators
$a(1+w_j+w_k)$ have the same residue modulo $2^r$,
since $k\ge r$ and $a$ is odd. Their residues modulo $M$
are distinct. In centered representatives, the condition
$\|y_{jk}\|<2^{-u}$ places them in the interval
$(-2^r,2^r)$, possibly intersected with the centered
interval of length $M$. Each residue class modulo $2^r$
has at most three representatives in that interval.
Thus there are at most $3e$ good mixed pairs.
All late--late pair numerators are congruent to $a$
modulo $2^r$, and distinct by the subset-sum injectivity.
There are at most three good late--late pairs. Consequently
\[
             N_u\le\binom r2+3r+3.
 \tag{P-F57f4}
\]

There is a loss when the anchor is separated from a cube
root. Set $\tau=\|3\theta\|=J/M$ and suppose
$\tau>3\cdot2^{-u}$.
Partition the $e$ early indices into
$h=\lfloor e/2\rfloor$ disjoint consecutive pairs,
starting with $(2,3)$, and discard at most one final index.
For each two distinct index-pairs $(j,j+1)$ and $(k,k+1)$,
the corresponding four pair factors are distinct and
belong to (P-F57f3). The identity
\[
          y_{jk}+y_{j+1,k}+y_{j,k+1}=3\theta
                                      \pmod1
 \tag{P-F57f5}
\]
follows from $w_{j+1}=-2w_j$ and $w_{k+1}=-2w_k$.
The three displayed phases cannot all have distance
less than $2^{-u}$ from an integer. Thus at least one
of the four factors in each block is bad.
The blocks for distinct unordered pairs of index-pairs
are disjoint as sets of unordered factor pairs.
There are exactly $\binom h2$ such blocks. Therefore,
under this anchor condition, (P-F57f4) improves to
\[
       N_u\le\binom r2+3r+3
              -\binom{\lfloor(r-2)_+/2\rfloor}{2}.
 \tag{P-F57f6}
\]
The omitted factor $(2,3)$ is internal to one index-pair
and is not used in any of these four-cell blocks.

For any nonzero dyadic phase $y$ with denominator dividing
$M$, one has $\|y\|\ge1/M$ and
\[
 \log_2|1-e^{2\pi iy}|^{-1}
 \le-\log_2\|y\|
 \le\sum_{u=0}^{l-1}\mathbf1_{\{\|y\|<2^{-u}\}}.
 \tag{P-F57f7}
\]
The last inequality is the elementary layer-cake bound:
the number of integers $u$ with $0\le u< -\log_2\|y\|$
is its ceiling. The right side is nonnegative, so it
also bounds the sum over the increment's subset of active
pairs after adding back its one omitted factor.
All other increment factors number $2n-4$.
Each is nonzero modulo $M$, and hence has logarithmic
reciprocal at most $l$. Their total contribution is
less than $2n^2$.

Put $K=(\lfloor\log_2J\rfloor-2)_+$.
For every integer $1\le r\le K$ one has
$J>3\cdot2^r$, so the anchor condition of (P-F57f6)
holds at $u=l-r$.
Summing (P-F57f4)--(P-F57f7) gives
\[
 \log_2 I_n(a)
 \le\binom n3+\frac32l(l+1)+3l+2n^2
       -\sum_{r=1}^K
                  \binom{\lfloor(r-2)_+/2\rfloor}{2}.
 \tag{P-F57f8}
\]
The first four terms are at most $n^3/6+4n^2$ for $n\ge5$.
For $r\ge6$, the summand is at least $(r-5)^2/8$:
$\lfloor(r-2)/2\rfloor\ge(r-3)/2$, and the factors
in $h(h-1)/2$ are nonnegative in this range.
Consequently the sum is at least
$(K-5)_+^3/24\ge(\log_2J-8)_+^3/24$.

It remains to compare this upper bound with the selected
increment. Among its factors, the
$B=\binom{n-3}{2}-1$ active pair weights in (P-F57f3),
excluding $(2,3)$, are divisible by three.
At $a=b_n$ the factor for largest index $k$ is at most
$8\cdot2^k/M$, exactly as in (P-F56ah).
The sum of $n-1-k$ over these pairs is
$\binom{n-2}{3}-(n-4)$.
Every remaining increment factor has absolute value at
most two. There are $\binom{n-1}{2}$ factors in total.
It follows that
\[
 \begin{split}
 \log_2 I_n(b_n)
 &\ge\binom{n-2}{3}-(n-4)-2B-\binom{n-1}{2}\\
 &=n^3/6-3n^2+71n/6-11
 \ge n^3/6-3n^2.
 \end{split}
 \tag{P-F57f9}
\]
Subtracting proves (P-F57f1).
If $\log_2J\ge\delta n$ and $\delta n\ge16$,
then $(\log_2J-8)^3/24\ge\delta^3n^3/192$.
If also $n\ge2688/\delta^3$, then
$7n^2\le\delta^3n^3/384$.
These two inequalities prove (P-F57f2), with its explicit
sufficient dimension threshold.\hfill$\square$
